\documentclass[a4paper]{article}
\usepackage[margin=1in]{geometry} % 设置边距，符合Word设定
\usepackage{ctex}
\usepackage{lipsum}
\usepackage{calc}
\usepackage{xcolor}
\usepackage{geometry}
\usepackage{tcolorbox}
\pagestyle{empty}
\pagenumbering{gobble} %此行可以取消页码

\title{\heiti\zihao{2} 机器人学习}
\author{\songti coco}
\date{4/3/2026}
\begin{document}
    \maketitle
\section{机器人}

% \newcommand{\tnss}{The not so Short Introduction to \LaTeXe}

% This is ``\tnss'' \ldots{} ``\tnss''.

% \newcommand{\note}{\noindent\textbf{注：}}
% \note{This is a text!}

% \newcommand{\degree}[1]{#1$^\circ$}
% \degree{30}

\newcommand{\textred}[1]{\textcolor{red}{#1}\par}
%\textred{这一行字体会变成红色。}

% \newenvironment{redenv}{
%     \color{red}
% }{
%     \par
% }
% % 红色字体环境

% \begin{redenv}
% 这是一段红色的字体。
% \end{redenv}

机器人操作系统（Robot Operating System，ROS）是一个灵活的软件
框架，为机器人软件开发提供了丰富的工具和软件库，可以帮助人们提高开
发机器人的效率。

选择ROS的一些主要原因（高端功能）：ROS自带现成的算法。

比如，即时定位与地图构建
（Simultaneous Loc-ali-zation and Mapping，SLAM）和自适应蒙特卡罗定
位（Adaptive Monte Carlo Localiza-tion，AMCL），这两个软件包可用于
移动机器人的自主导航。

ROS的消息传递机制允许在不同节点间进行通信。我们可以用
C++或C语言对节点进行编程，当然也可以用Python或Java语言。

\textred{ROS技术社区（http://answers.ros.org）}

ROS的主要仿真平台是Gazebo，Gazebo仿真器没有内置的编程
环境，所有仿真必须在ROS下编程完成。与Gazebo比较，其他仿真环境（如
V-REP、Webots）有内置的编程环境，可以直接对机器人进行原型设计与编
程。而且，这些仿真环境有丰富的图形界面（GUI）工具集，支持各类机器
人，并且提供了ROS接口。虽然这些相应的工具都是各自平台专用的，但效
果也都不错。
ROS中主要有三个级别：文件系统、计算图、社区。

为了能创建、修改、使用ROS软件包，我们需要先了解一些常用命令。
下面是一些使用ROS软件包的常用命令
catkin\_create\_pkg 此命令用于创建一个新软件包。
rospack 此命令用于获取软件包相关的信息。
catkin\_make 此命令用于在工作区中编译软件包。
rosdep 此命令将安装一个软件包所需的系统依赖项。

ROS博客会定期更新与ROS社区相关的新闻、照片、视频
（http://www.ros.org/news）

\section{ROS编程入门}
ROS软件包是ROS系统的基本单位。

使用ROS软件包的第一个要求是创建ROS的catkin工作区。

话题是两个节点间通信的基本方式。
下面我们将创建两个ROS节点，一个发布话题，另一个订阅该话题。

节点间的相互通讯通过/number完成。
demo\_topic\_publisher节点向/numbers话题发布信息，然后
demo\_topic\_subscriber节点订阅这个话题

与ROS服务一样，在actionlib中，我们也必须初始化动作规范。这个动
作的规范存储在以.action结尾的文件中。该文件必须保存在ROS软件包内的
action文件夹中。action文件包含以下各部分内容。
·Goal（目标）：动作客户端可以发送一个必须由动作服务器来执行的目标。
这就类似于ROS服务中的请求。例如，如果机器人手臂关节想从45度转动到
90度，那么这里的目标就是90度。
·Feedback（反馈）：动作客户端向动作服务器发送目标后，将开始执行回
调函数。反馈只是简单地给出回调函数内当前操作的进度。通过使用反馈，
我们可以获得当前任务的进度。在前面的例子中，机器人手臂关节必须移动
到90度，在这种情况下，反馈就是机械臂从45度转动到90度之间的中间
值。
·Result（结果）：完成目标后，动作服务器将发送完成的最终结果，它可以
是计算结果或者一个确认。在前面的例子中，如果关节转动到90度，则完成
了目标任务，结果就是可以表示目标完成的任何形式。

话题（topic）、服务（service）和动作库（actionlib）可应用于不同的
场景。我们知道话题是单向通信的，服务是具备请求/应答功能的双向通
信，动作库是ROS服务的一种变化形式，我们可以在需要时取消在服务器上
运行的进程。

\section{在ROS中为3D机器人建模}
机器人制造的第一个阶段是设计和建模。这里我们可以使用CAD软件对一个机器人进行设
计和建模，例如AutoCAD、SOLIDWORKS和Blender。机器人建模的主要目的之一就是仿
真。

将讨论两类机器人的设计。一个是七 自由度（7 Degrees of Freedom，7-
DOF）机械手臂，另一个是差速驱动机器人。

可以使用URDF来定义机器人模型、传感器和工作环境，使用URDF
解析器对其进行解析。

设计好URDF后，可以在RViz上查看它。我们可以创建一个
view\_demo.launch启动文件，并将下面的代码存放到该文件中，然后将该文
件放入launch文件夹。

\end{document}